%============================================================
%  Chapter 10 – Security and Authentication
%============================================================
\chapter{Security and Authentication}

\section{Overview}

Security was a primary design consideration throughout the development of \projectname{}. This chapter documents all authentication mechanisms, authorisation controls, and security-relevant implementation decisions.

\section{Password Security}

All user passwords are hashed using \textbf{bcrypt.js} \cite{bcryptjs} before storage. bcrypt incorporates a configurable work factor (cost factor), which is set to \textbf{12} in \projectname{}, making brute-force attacks computationally expensive. The library automatically generates and embeds a cryptographically random salt in each hash, preventing rainbow table attacks.

\begin{lstlisting}[language=TypeScript, caption={Password Hashing During Signup}]
const SALT_ROUNDS = 12;
const passwordHash = await bcrypt.hash(plainTextPassword, SALT_ROUNDS);
\end{lstlisting}

Plain-text passwords are never stored, logged, or transmitted after the initial POST request body.

\section{JWT Authentication}

\projectname{} uses a \textbf{dual-token} JWT authentication scheme:

\begin{table}[H]
\centering
\caption{JWT Token Properties}
\begin{tabularx}{\textwidth}{llX}
\toprule
\textbf{Token Type} & \textbf{Lifetime} & \textbf{Storage Location} \\
\midrule
Access Token  & 15 minutes & JavaScript memory (React context). Never \texttt{localStorage}. \\
Refresh Token & 7 days     & HttpOnly cookie (inaccessible to JavaScript; sent automatically by the browser). \\
\bottomrule
\end{tabularx}
\end{table}

Both tokens are signed with separate secrets (\texttt{JWT\_ACCESS\_SECRET} and \texttt{JWT\_REFRESH\_SECRET}) stored as environment variables.

\subsection{Token Rotation}

On each call to \texttt{POST /api/auth/refresh}, the backend:

\begin{enumerate}
  \item Reads the refresh token from the HttpOnly cookie.
  \item Looks up the corresponding \texttt{Session} document.
  \item If found and not expired: deletes the old session, generates a new refresh token, creates a new \texttt{Session}, sets a new HttpOnly cookie, and returns a new access token.
  \item If the refresh token is reused (already deleted from the DB), a replay attack is detected and all sessions for the user are revoked immediately.
\end{enumerate}

\section{Google OAuth 2.0}

Google authentication is implemented using \textbf{Passport.js} \cite{passport2024} with the \texttt{passport-google-oauth20} strategy \cite{oauth2}. The flow proceeds as follows:

\begin{enumerate}
  \item User clicks ``Continue with Google'' on the frontend.
  \item Browser is redirected to \texttt{GET /api/auth/google}, which triggers the Passport redirect to Google's authentication servers.
  \item User authenticates with Google; Google redirects back to \texttt{GET /api/auth/google/callback} with an authorisation code.
  \item Passport exchanges the code for a Google profile, using the \texttt{googleId} to find or create a \texttt{User} document.
  \item JWT tokens are generated and the browser is redirected to \texttt{/auth/google/success?token=...} with the access token as a URL parameter for a brief handshake (then removed from URL).
\end{enumerate}

The client credentials (\texttt{GOOGLE\_CLIENT\_ID} and \texttt{GOOGLE\_CLIENT\_SECRET}) are stored as environment variables only.

\section{OTP Email Verification}

Upon registration, a 6-digit OTP is generated using \texttt{crypto.randomInt(100000, 999999)}. It is:

\begin{itemize}
  \item Stored in Redis with a 10-minute TTL under the key \texttt{otp:\{email\}}.
  \item Sent via Nodemailer to the user's email as an HTML-formatted message.
  \item Deleted from Redis upon successful verification.
\end{itemize}

Failed OTP attempts are not currently rate-limited (planned future enhancement), but the short TTL limits the attack window.

\section{Role-Based Access Control (RBAC)}

All protected API routes are guarded by two layers of middleware:

\begin{enumerate}
  \item \textbf{\texttt{protect}} — verifies the JWT access token and attaches the decoded \texttt{\{userId, role\}} to \texttt{req.user}.
  \item \textbf{\texttt{requireRole(role)}} — checks that \texttt{req.user.role === role}; returns 403 if not.
\end{enumerate}

The frontend enforces the same constraints via \texttt{RoleRoute} components that render a 403 screen or redirect if the user's role does not match.

\section{Environment Variable Security}

All secrets and configuration values are stored in a \texttt{.env} file:

\begin{itemize}
  \item \texttt{.env} is listed in \texttt{.gitignore} to prevent accidental commitment.
  \item Environment variables are accessed exclusively via \texttt{process.env.*} after calling \texttt{dotenv.config()}.
  \item The \texttt{.env} file is never exposed to the frontend (Vite/React receives only \texttt{VITE\_*} prefixed variables, and none of the backend secrets carry this prefix).
\end{itemize}

See Appendix~\ref{app:env} for a full list of required environment variables.

\section{CORS Configuration}

Cross-Origin Resource Sharing is configured to allow requests only from the specific frontend origins. The current development configuration allows:

\begin{itemize}
  \item \texttt{http://localhost:5173} (Vite dev server)
  \item \texttt{http://localhost:3000}
  \item \texttt{http://localhost:8080}
\end{itemize}

\texttt{credentials: true} is enabled to allow the browser to send HttpOnly cookies with cross-origin requests.

\section{File Upload Security}

\begin{itemize}
  \item Multer enforces a \textbf{10MB} maximum file size.
  \item Files are stored in Cloudinary, not on the application server. The server never stores files on its filesystem.
  \item Uploaded file types are not explicitly validated in the current version (planned future enhancement).
\end{itemize}

\section{XSS and Injection Prevention}

\begin{itemize}
  \item The access token is stored in JavaScript memory, not \texttt{localStorage}, preventing XSS token theft.
  \item All request body data is validated via \textbf{Zod} schemas before database operations.
  \item Prisma's parameterised query model prevents SQL/NoSQL injection at the ORM level.
\end{itemize}

\section{Account Security Features}

\begin{itemize}
  \item \textbf{Account ban}: Admins can ban accounts. Banned users receive a 403 response on login with the ban reason.
  \item \textbf{ID verification}: A boolean \texttt{isIdVerified} field is set by admins after reviewing the uploaded ID document.
  \item \textbf{Password reset}: A UUID-based reset token (stored hashed in \texttt{PasswordResetToken}) with a 1-hour expiry. Tokens are single-use (marked \texttt{used: true} on consumption).
\end{itemize}
